\section{Data Augmentation}
\label{sec:data}

The SoccerNet Player-Centric Ball Action Spotting (PCBAS) challenge~\cite{deliege2021soccernet,cioppa2024soccernet} extends the standard action spotting task by requiring systems to identify not only \emph{what} ball-related action occurs and \emph{when}, but also \emph{who} performs it. Eight action classes are evaluated---drive, pass, cross, throw-in, shot, header, tackle, block---spanning highly frequent events (pass: 3059 train examples) to rare ones (tackle: 26). The provided tactical data gives 13-column per-frame tracking at 1920$\times$1080, 25fps, but $\sim$22\% of tracked players are missing bounding boxes and no ball position is provided. We address both gaps with SAM3-based~\cite{ravi2024sam2} pipelines.

\noindent\textbf{Bounding box recovery.}
For frames with $\geq$4 known bboxes, we compute a foot-point homography $H$ mapping pitch coordinates to pixel space at the bottom-center of each bbox. Missing players are projected through $H$ and classified as \emph{out-of-frame} (outside 1920$\times$1080 $\pm$ 50px), \emph{replay} ($<$4 reference bboxes), or \emph{in-frame}. In-frame candidates are matched to SAM3 ``person'' detections (conf $\geq 0.3$) filtered by size (area $\geq 80$px, $h/w \geq 0.7$, $w \leq 250$, $h \leq 400$) and cross-validated against existing GT bboxes (IoU $> 0.2$). Remaining detections are assigned by foot-point Euclidean distance ($\leq 150$px). SAM3 runs every 5th frame with $\pm 2$-frame propagation, filling $\sim$6.8\% of missing rows---the genuinely visible but untracked players.

\noindent\textbf{Ball detection.}
SAM3 with ``soccer ball'' (conf $\geq 0.15$) generates per-frame candidates filtered by size ($\leq$50px), aspect ratio ($\leq$2.5), dynamic pitch bounds from player positions, and proximity to nearest player ($\leq$300px). Viterbi DP selects the optimal trajectory: emission costs combine detection confidence and a player-proximity bonus (up to 40 points for $<$80px); transition costs penalize teleportation ($0.3\times$ for $\leq$120px/frame, $1\times$ for $\leq$360px, $5\times$ beyond). A 5-frame median filter and linear interpolation across gaps $\leq$4 frames smooth the output.

\section{Feature Engineering}
\label{sec:features}

\noindent\textbf{Motion (34-dim).} From raw per-player tracking $(x, y, v_x, v_y)$ we compute: \emph{kinematics} (13-dim: position, velocity, speed, acceleration, jerk, movement direction $\theta = \text{atan2}(v_y, v_x)/\pi$, direction change), \emph{identity} (14-dim: 13-role one-hot + team indicator), and \emph{inter-player} (7-dim: nearest teammate/opponent distances, relative velocity to nearest opponent, opponent pressure count within radius $< 0.05$, attacking and defending goal proximity).

\noindent\textbf{Visual (768-dim).} Frozen DINOv2-ViT-B/14~\cite{oquab2024dinov2} extracts patch features at 518$\times$518 resolution. Per-player ROI-Align ($3{\times}3$) over bounding boxes produces 768-dim embeddings. A binary visibility mask is appended (769-dim input). For players with missing bboxes, a learned 768-dim \texttt{missing\_vis} embedding replaces the zero-vector.

\noindent\textbf{Ball (6-dim).} Per-player ball features from detected positions (Sec.~\ref{sec:data}): distance to ball, relative displacement $\Delta x/\Delta y$ normalized by image dimensions, binary \texttt{is\_closest} flag, and ball velocity $(v_x, v_y)$ via $\pm$3-frame finite differences. This was the single largest additive feature (+5.5 F1 points).

\noindent\textbf{Scene (768-dim).} The DINOv2 CLS token provides a global scene descriptor capturing pitch region, camera zoom, and broadcast state. This is broadcast identically to all 22 player slots before fusion.

\section{Model Architecture}
\label{sec:model}

\noindent\textbf{Stage 1: Event detection.}
A DeepSets~\cite{zaheer2017deep} + dilated temporal conv network operates on 41-dim features (the 34-dim motion features plus goalkeeper flag, team centroid, and padding). Per-player MLP embeddings are aggregated via max+mean pooling, then 10 dilated 1D residual blocks~\cite{he2016deep} (GroupNorm, GELU, dilations $2^0$--$2^9$, receptive field $\sim$1024 frames) produce per-frame event scores. Candidates: $\sigma(\text{score}) > 0.3$, de-duplicated to one per 12 frames.

\noindent\textbf{Stage 2: Multi-modal classification.}
Per candidate, a 100-frame window feeds four branches: motion (MLP $34 \to 128$ + 6 BiBlocks, center-frame extraction), visual ($769 \to 128$), ball ($6 \to 32$), scene ($768 \to 128$). All branches are concatenated per player and projected to 256-dim, then two transformer encoder layers~\cite{vaswani2017attention} (4 heads, FFN 512, dropout 0.15) model cross-player interactions. Action class is predicted from max-pooled features; player identity from per-slot logits. Deep variants use 384-dim, 8 heads, and 3 transformer layers.

\noindent\textbf{Robustness to missing visual.}
We train robust variants with: (1) a learned 768-dim \texttt{missing\_vis} embedding replacing zero-vectors for missing bboxes, (2) visual dropout---25\% of batches replace all visual with \texttt{missing\_vis}, simulating no-visual frames, and (3) KL consistency regularization between predictions with/without visual. The \textbf{RobustScene} family combines all three with the scene branch.

\noindent\textbf{Training.}
All models are trained on train+val combined (102 game halves, $\sim$6200 events). We use cross-entropy with class weights $[0.1, 1, 0.8, 5, 5, 8, 3, 20, 8]$ reflecting severe class imbalance---pass has 3059 examples while tackle has only 26. AdamW~\cite{loshchilov2019adamw} (lr=$5{\times}10^{-4}$, wd=$10^{-4}$), batch size 64, 25 epochs. Temporal jitter ($\pm$3 frames) augments event frame alignment. We train 25 model variants spanning 8 architecture families (base, deep, wide, ball-window, ball-v2, GSF, robust, scene-robust) across multiple seeds ($\{42, 123, 456, 789\}$), yielding diverse predictions for ensembling.

\noindent\textbf{Inference.}
Per-class decoding with optimized thresholds. TTA averages softmax over 5 offsets $\{-4,-2,0,+2,+4\}$. Per-class ensemble selects 2--12 of 25 variants via greedy leave-one-out. RobustScene improves most classes but degrades tackle precision (0.40 $\to$ 0.20), so it is excluded from that subset---ensemble composition must be tuned per-class. Per-class NMS (12-frame window) removes duplicates.

\noindent\textbf{Results} are shown in Table~\ref{tab:challenge}. Recall is strongly affected by bbox availability (drive: 0.761 with bbox vs.\ 0.401 without; throw-in: 0.913 vs.\ 0.755) and replay frames (drive: 0.497, header: 0.167), motivating our bbox-filling and robustness pipelines. The micro-macro F1 gap (0.712 vs.\ 0.557) reflects class imbalance.

Key findings from 200+ experiments: (1) tracking data provides 80\%+ of signal---a motion-only model already reaches F1 $\sim$0.63; (2) ball proximity is the single largest additive feature (+5.5 F1), while more complex ball encodings (IoU overlap, altitude, 8-dim variants) offered no improvement over the simple 6-dim distance-based features; (3) architectural diversity (different branch designs, depths, feature variants) outperforms seed diversity for ensembles; (4) temporal video features (VideoMAE~\cite{tong2022videomae}, ViViT, V-JEPA2~\cite{assran2025vjepa2}) consistently failed to improve over single-frame DINOv2, likely because tracking data already captures temporal dynamics; (5) alternative image backbones (DINOv3, SigLIP2, CRADIO) all performed comparably to DINOv2-B---backbone choice does not matter when ROI features are pooled to a single vector; (6) rare classes (tackle: 26 examples, block: 78) are fundamentally data-limited---no training trick (focal loss, asymmetric loss, oversampling, curriculum learning, hard negative mining, contrastive losses) compensated for the scarcity; (7) post-processing heuristics (transition constraints, position-based gating, cross-class voting) either hurt rare classes or showed no effect, confirming that the model's learned representations already capture the relevant patterns.

\begin{table}[h]
\centering
\small
\begin{tabular}{lccc}
\hline
\textbf{Class} & \textbf{P} & \textbf{R} & \textbf{F1} \\
\hline
Drive & .682 & .704 & .693 \\
Pass & .809 & .739 & .773 \\
Cross & .850 & .788 & .818 \\
Throw-in & .766 & .828 & .796 \\
Shot & .606 & .597 & .602 \\
Header & .382 & .448 & .412 \\
Tackle & .400 & .083 & .138 \\
Block & .213 & .235 & .224 \\
\hline
\textbf{Macro} & .589 & .553 & \textbf{.557} \\
\textbf{Micro} & .721 & .703 & .712 \\
\hline
\end{tabular}
\caption{Challenge set results (6 game halves).}
\label{tab:challenge}
\end{table}
